\chapter{بحث}
\label{ch:discussion}

\section{یافته‌های اصلی}
این پایان‌نامه نشان داد که در کوهورت کامل \lr{MIMIC\text{-}IV}، داده‌های \pnum{24} ساعت نخست \lr{ICU} می‌توانند نشانگر جانشین امکان‌پذیری ترخیص را با عملکرد قابل‌توجه اما نه قطعی پیش‌بینی کنند. بالاترین \lr{AUROC} متعلق به مدل انباشتی کلاسیک بود، اما \lr{XGBoost} از نظر کالیبراسیون و امتیاز بریر قابل‌دفاع‌تر بود. این تمایز مهم است، زیرا در پژوهش‌های سلامت، احتمال خروجی مدل تنها زمانی مفید است که تا حد امکان با خطر مشاهده‌شده هم‌خوان باشد.

\section{چرا \lr{XGBoost} مدل اصلی قابل‌دفاع‌تر است؟}
مدل انباشتی با \lr{AUROC} \cnum{0.755} اندکی بالاتر از \lr{XGBoost} بود، اما \lr{ECE} آن \cnum{0.109} بود. در مقابل، \lr{XGBoost} با \lr{AUROC} \cnum{0.753} و \lr{ECE} \cnum{0.020} توازن بهتری ارائه کرد. در یک سناریوی پژوهشی نزدیک به پشتیبان تصمیم، چنین توازنی مهم‌تر از بهبود بسیار کوچک \lr{AUROC} است. به همین دلیل، این پایان‌نامه \lr{XGBoost} را مدل اصلی قابل‌دفاع‌تر می‌داند، نه الزاماً مدل با بیشترین عدد در یک معیار واحد.

\section{چرا \lr{Few-Shot ETHOS} ضعیف‌تر بود؟}
نتیجۀ ضعیف‌تر \lr{Few-Shot ETHOS} یافته‌ای منفی اما ارزشمند است. یادگیری چندنمونه‌ای در حوزه‌هایی موفق است که بازنمایی‌های پرظرفیت و ساختارهای شباهت پیچیده سودمند باشند. در این وظیفه، داده‌ها جدولی، کم‌بعد و دارای معنای مستقیم بالینی بودند. مدل‌های درختی می‌توانستند آستانه‌ها، تعامل‌ها و غیرخطی‌ها را به‌خوبی استخراج کنند. از سوی دیگر، توکن‌سازی فشرده ممکن است بخشی از شدت و پیوستگی مقادیر آزمایشگاهی را کاهش داده باشد. بنابراین، \lr{ETHOS} خام نتوانست با مدل‌های جدولی رقابت کند.

این نتیجه به‌معنای بی‌ارزش بودن توکن‌ها نیست. عملکرد \lr{Token Hybrid RF} نشان داد وقتی توکن‌ها در کنار ویژگی‌های جدولی استفاده شوند، می‌توانند سیگنال مکمل داشته باشند. مسیر آینده باید بر آموزش نظارت‌شدۀ بهتر توکن‌ها، اپیزودهای بالینی معنادارتر، و هم‌جوشی کالیبراسیون‌محور تمرکز کند.

\section{پیامدهای مدل هیبریدی و انباشتی}
مدل هیبریدی و مدل انباشتی نشان دادند ترکیب بازنمایی‌ها می‌تواند مفید باشد. بااین‌حال، بهبود مدل انباشتی باید با احتیاط تفسیر شود؛ زیرا کالیبراسیون آن از \lr{XGBoost} ضعیف‌تر بود. در کاربردهای پژوهشی، مدل انباشتی می‌تواند برای رتبه‌بندی خطر مناسب باشد، اما برای گزارش احتمال، نیازمند کالیبراسیون جداگانه است.

\section{یادگیری فدرال}
شبیه‌سازی فدرال نشان داد مدل گره‌ای لجستیک می‌تواند به رگرسیون لجستیک متمرکز نزدیک بماند. این یافته برای حریم خصوصی مهم است، زیرا نشان می‌دهد اگر هدف یک مدل ساده و قابل‌توضیح باشد، توزیع گره‌ای الزاماً افت شدید ایجاد نمی‌کند. اما این نتیجه نباید به استقرار واقعی تعمیم داده شود. در بیمارستان‌های واقعی، گره‌ها ناهمگن‌ترند، اندازه نمونه‌ها متفاوت است، ارتباطات محدود است و نیاز به ارزیابی امنیت، انصاف و پایش انحراف وجود دارد.

\section{رابطه با عنوان پایان‌نامه}
عنوان پایان‌نامه از «فیوشات درمانی» و «قابلیت درمان» سخن می‌گوید. بازنگری علمی متن لازم داشت که این عنوان با اجرای واقعی همسو شود. بنابراین، «قابلیت درمان» در این پایان‌نامه به‌صورت «نشانگر جانشین امکان‌پذیری ترخیص از \lr{ICU}» تعریف شده است. این شفاف‌سازی از بیش‌ادعایی جلوگیری می‌کند و نشان می‌دهد پژوهش یک سامانه درمان‌گر یا توصیه‌گر درمان نیست، بلکه یک مطالعۀ مدل‌سازی گذشته‌نگر است.

\section{نقاط قوت}
نقاط قوت پژوهش شامل استفاده از کوهورت کامل واجد شرایط، حذف سقف نمونه‌گیری، استفاده از داده خام کامل و مانیفست‌تأییدشده، مقایسۀ چند خانوادۀ مدل، گزارش کالیبراسیون، و تفکیک روشن میان نتایج اصلی و بنچمارک تاریخی است. همچنین متن حاضر نتایج منفی را پنهان نکرده و نشان داده است که مدل‌های کلاسیک در داده‌های جدولی بالینی همچنان باید جدی گرفته شوند.

\section{محدودیت‌ها}
محدودیت‌های اصلی عبارت‌اند از:
\begin{enumerate}
\item مطالعه گذشته‌نگر و مبتنی بر یک پایگاه داده است.
\item برچسب، جانشین است و توسط متخصصان به‌عنوان آمادگی ترخیص یا امکان‌پذیری واقعی درمان داوری نشده است.
\item اعتبارسنجی خارجی روی بیمارستان دیگر انجام نشده است.
\item اعتبارسنجی آینده‌نگر یا مداخله‌ای وجود ندارد.
\item ویژگی‌های \pnum{24} ساعت نخست مسیرهای دیرتر بیماری را کامل نشان نمی‌دهند.
\item احتمال مخدوش‌سازی باقیمانده وجود دارد.
\item یادگیری فدرال شبیه‌سازی شده و در شبکه واقعی بیمارستانی اجرا نشده است.
\item \lr{Few-Shot ETHOS} در این نسخه ضعیف‌تر از مدل‌های جدولی بود و نیاز به بهبود دارد.
\item کالیبراسیون میان مدل‌ها متفاوت است و \lr{AUROC} به‌تنهایی کافی نیست.
\item دادۀ خام \lr{MIMIC\text{-}IV} طبق محدودیت‌های استفاده قابل بازتوزیع نیست.
\end{enumerate}

\section{کارهای آینده}
مسیرهای آینده شامل اعتبارسنجی خارجی، اعتبارسنجی آینده‌نگر، مدل‌سازی سری زمانی غنی‌تر، آموزش نظارت‌شدۀ توکن‌ها، ساخت اپیزودهای بالینی معنادارتر برای یادگیری چندنمونه‌ای، آموزش کالیبراسیون‌محور، تحلیل انصاف و زیرگروه، فدرال واقعی چندمرکزی، و ارزیابی انسان-در-حلقه با بالین‌گران است.
